% Source: physical pp. 263--270 (printed pp. 240--247). Coverage:
% Problems 1--5 and References 1--45, including all lettered entries,
% ending before Chapter 29. Figures/tables/footnotes/dividers: none.

\chapterexercisehook{28}

\chapterbackmatter{References}

\begin{enumerate}
  \item \phantomsection\label{ch28-ref-1}
  We will not consider here the possibility of unification at much
  lower energies raised by I.~Antoniadis, \textit{Phys. Lett.}
  B\textbf{246}, 377 (1990); J.~Lykken, \textit{Phys. Rev.}
  D\textbf{54}, 3693 (1996), and revived by N.~Arkani-Hamed,
  S.~Dimopoulos, and G.~Dvali, \textit{Phys. Lett.} B\textbf{429}, 263
  (1998); K.~R.~Dienes, E.~Dudea, and T.~Gherghetta,
  \textit{Phys. Lett.} B\textbf{436}, 55 (1998); I.~Antoniadis,
  N.~Arkani-Hamed, S.~Dimopoulos, and G.~Dvali,
  \textit{Phys. Rev. Lett.} B\textbf{436}, 257 (1998).

  \item[1a.] \phantomsection\label{ch28-ref-1a}
  S.~Weinberg, in \textit{Proceedings of the XVII International
  Conference on High Energy Physics, London, 1974}, J.~R.~Smith, ed.
  (Rutherford Laboratory, Chilton, Didcot, England, 1974); S.~Weinberg,
  \textit{Phys. Rev.} D\textbf{13}, 974 (1976); E.~Gildener and
  S.~Weinberg, \textit{Phys. Rev.} D\textbf{13}, 3333 (1976).

  \item[1b.] \phantomsection\label{ch28-ref-1b}
  T.~Banks and L.~Dixon, \textit{Nucl. Phys.} B\textbf{307}, 93
  (1988). For a detailed treatment, see J.~Polchinski,
  \textit{String Theory} (Cambridge University Press, Cambridge,
  1998): Chapter 18.

  \item \phantomsection\label{ch28-ref-2}
  Additive $R$-conservation laws were introduced by A.~Salam and
  J.~Strathdee, \textit{Nucl. Phys.} B\textbf{87}, 85 (1975);
  P.~Fayet, \textit{Nucl. Phys.} B\textbf{90}, 104 (1975), reprinted
  in \textit{Supersymmetry}, S.~Ferrara, ed., (North Holland/World
  Scientific, Amsterdam/Singapore, 1987). The $R$ parity can be
  defined in terms of an $R$ quantum number as $\exp(\ii\pi R)$, and may
  be conserved multiplicatively even if $R$ is not conserved
  additively; see G.~Farrar and P.~Fayet, \textit{Phys. Lett.}
  \textbf{76B}, 575 (1978); P.~Fayet, in \textit{Unification of the
  Fundamental Particle Interactions}, S.~Ferrara, J.~Ellis, and
  P.~van Nieuwenhuizen, eds. (Plenum, New York, 1980); S.~Dimopoulos,
  S.~Raby, and F.~Wilczek, \textit{Phys. Lett.} \textbf{112B}, 133
  (1982); G.~Farrar and S.~Weinberg, \textit{Phys. Rev.}
  D\textbf{27}, 1731 (1983), reprinted in \textit{Supersymmetry}.

  \item \phantomsection\label{ch28-ref-3}
  S.~Dimopoulos and H.~Georgi, \textit{Nucl. Phys.} B\textbf{193}, 150
  (1981), reprinted in \textit{Supersymmetry},
  Ref.~\hyperref[ch28-ref-2]{2}.

  \item \phantomsection\label{ch28-ref-4}
  R.~D.~Peccei and H.~Quinn, \textit{Phys. Rev. Lett.} \textbf{38},
  1440 (1977); \textit{Phys. Rev.} D\textbf{16}, 1791 (1977).

  \item[4a.] \phantomsection\label{ch28-ref-4a}
  S.~Dimopoulos and G.~F.~Giudice, \textit{Phys. Lett.} B\textbf{357},
  573 (1995); A.~Pomerol and D.~Tommasini, \textit{Nucl. Phys.}
  B\textbf{466}, 3 (1996); G.~Dvali and A.~Pomerol,
  \textit{Phys. Rev. Lett.} \textbf{77}, 3728 (1996);
  \textit{Nucl. Phys.} B\textbf{522}, 3 (1998); A.~G.~Cohen,
  D.~B.~Kaplan, and A.~E.~Nelson, \textit{Phys. Lett.} B\textbf{388},
  588 (1996); R.~N.~Mohapatra and A.~Riotto, \textit{Phys. Rev.}
  D\textbf{55}, 1 (1997); R.-J.~Zhang, \textit{Phys. Lett.}
  B\textbf{402}, 101 (1997); H-P.~Nilles and N.~Polonsky,
  \textit{Phys. Lett.} B\textbf{412}, 69 (1997); D.~E.~Kaplan,
  F.~Lepeintre, A.~Masiero, A.~E.~Nelson, and A.~Riotto,
  hep-ph/9806430, to be published; J.~Hisano, K.~Kurosawa, and
  Y.~Nomura, \textit{Phys. Lett.} B\textbf{445}, 316 (1999). This
  pattern of masses can arise naturally from radiative corrections;
  see J.~L.~Feng, C.~Kolda, and N.~Polonsky, \textit{Nucl. Phys.}
  B\textbf{546}, 3 (1999); J.~Bagger, J.~L.~Feng, and N.~Polonsky,
  hep-ph/9905292, to be published.

  \item[4b.] \phantomsection\label{ch28-ref-4b}
  B.~W.~Lee and S.~Weinberg, \textit{Phys. Rev. Lett.} \textbf{39},
  165 (1977); D.~A.~Dicus, E.~W.~Kolb, and V.~L.~Teplitz,
  \textit{Phys. Rev. Lett.} \textbf{39}, 168 (1977).

  \item[4c.] \phantomsection\label{ch28-ref-4c}
  S.~Wolfram, \textit{Phys. Lett.} \textbf{82B}, 65 (1979); J.~Ellis,
  J.~S.~Hagelin, D.~V.~Nanopoulos, K.~Olive, and M.~Srednicki,
  \textit{Nucl. Phys.} B\textbf{238}, 453 (1984).

  \item[4d.] \phantomsection\label{ch28-ref-4d}
  P.~F.~Smith and J.~R.~J.~Bennett, \textit{Nucl. Phys.} B\textbf{149},
  525 (1979).

  \item \phantomsection\label{ch28-ref-5}
  H.~Georgi, H.~R.~Quinn, and S.~Weinberg, \textit{Phys. Rev. Lett.}
  \textbf{33}, 451 (1974).

  \item \phantomsection\label{ch28-ref-6}
  S.~Dimopoulos and H.~Georgi, Ref.~\hyperref[ch28-ref-3]{3};
  J.~Ellis, S.~Kelley, and D.~V.~Nanopoulos, \textit{Phys. Lett.}
  B\textbf{260}, 131 (1991); U.~Amaldi, W.~de Boer, and H.~Furstmann,
  \textit{Phys. Lett.} B\textbf{260}, 447 (1991); C.~Giunti,
  C.~W.~Kim and U.~W.~Lee, \textit{Mod. Phys. Lett.} A\textbf{6}, 1745
  (1991); P.~Langacker and M.-X.~Luo, \textit{Phys. Rev.} D\textbf{44},
  817 (1991). For other references and more recent analyses of the data,
  see P.~Langacker and N.~Polonsky, \textit{Phys. Rev.} D\textbf{47},
  4028 (1993); D\textbf{49}, 1454 (1994); L.~J.~Hall and U.~Sarid,
  \textit{Phys. Rev. Lett.} \textbf{70}, 2673 (1993).

  \item \phantomsection\label{ch28-ref-7}
  S.~Dimopoulos, S.~Raby, and F.~Wilczek, \textit{Phys. Rev.}
  D\textbf{24}, 1681 (1981), reprinted in \textit{Supersymmetry},
  Ref.~\hyperref[ch28-ref-2]{2}.

  \item[7a.] \phantomsection\label{ch28-ref-7a}
  P.~Ho\v{r}ava and E.~Witten, \textit{Nucl. Phys.} B\textbf{460}, 506
  (1996); \textit{ibid.} B\textbf{475}, 94 (1996); E.~Witten,
  \textit{Nucl. Phys.} B\textbf{471}, 135 (1996); P.~Ho\v{r}ava,
  \textit{Phys. Rev.} D\textbf{54}, 7561 (1996).

  \item \phantomsection\label{ch28-ref-8}
  H.~Pagels and J.~R.~Primack, \textit{Phys. Rev. Lett.} \textbf{48},
  223 (1982).

  \item \phantomsection\label{ch28-ref-9}
  S.~Weinberg, \textit{Phys. Rev. Lett.} \textbf{48}, 1303 (1983).

  \item \phantomsection\label{ch28-ref-10}
  S.~Dimopoulos and H.~Georgi, Ref.~\hyperref[ch28-ref-3]{3};
  N.~Sakai, \textit{Z. Phys.} C\textbf{11}, 153 (1981). For reviews,
  see H.~E.~Haber and G.~L.~Kane, \textit{Phys. Reports} \textbf{117},
  75 (1985); J.~A.~Bagger, in \textit{QCD and Beyond: Proceedings of
  the Theoretical Advanced Study Institute in Elementary Particle
  Physics, University of Colorado, June 1995}, D.~E.~Soper, ed. (World
  Scientific, Singapore, 1996); V.~Barger, in \textit{Fundamental
  Particles and Interactions: Proceedings of the FCP Workshop on
  Fundamental Particles and Interactions, Vanderbilt University, May
  1997}, R.~S.~Panvini, T.~J.~Weiler, eds. (American Institute of
  Physics, Woodbury, NY, 1998); J.~F.~Gunion, in \textit{Quantum
  Effects in the MSSM---Proceedings of the International Workshop on
  Quantum Effects in the MSSM, Barcelona, September 1997}, J.~Sol\`a,
  ed. (World Scientific Publishing, Singapore, 1998); S.~Dawson, in
  \textit{Proceedings of the 1997 Theoretical Advanced Study Institute
  on Supersymmetry, Supergravity, and Supercolliders}, J.~Bagger, ed.
  (World Scientific, Singapore, 1998); S.~P.~Martin, in
  \textit{Perspectives on Supersymmetry}, G.~L.~Kane, ed. (World
  Scientific, Singapore, 1998); K.~R.~Dienes and C.~Kolda, in
  \textit{Perspectives on Supersymmetry}, \textit{ibid.}

  \item \phantomsection\label{ch28-ref-11}
  S.~Dimopoulos and D.~Sutter, \textit{Nucl. Phys.} B\textbf{194}, 65
  (1995); H.~Haber, \textit{Nucl. Phys. Proc. Suppl.} \textbf{62}, 469
  (1998).

  \item \phantomsection\label{ch28-ref-12}
  S.~Dimopoulos and H.~Georgi, Ref.~\hyperref[ch28-ref-3]{3};
  J.~Ellis and D.~V.~Nanopoulos, \textit{Phys. Lett.} \textbf{110B},
  44 (1982); J.~F.~Donoghue, H-P.~Nilles, and D.~Wyler,
  \textit{Phys. Lett.} \textbf{128B}, 55 (1983). For
  strong-interaction corrections to these calculations, see
  J.~A.~Bagger, K.~T.~Matchev, and R.-J.~Zhang,
  \textit{Phys. Lett.} B\textbf{412}, 77 (1997). Conditions under
  which limits on flavor changing processes do not constrain squark
  masses are discussed by R.~Barbieri and R.~Gatto,
  \textit{Phys. Lett.} \textbf{110B}, 211 (1981); Y.~Nir and
  N.~Seiberg, \textit{Phys. Lett.} B\textbf{309}. 337 (1993). For a
  detailed review, see F.~Gabbiani, E.~Gabrielli, A.~Masiero, and
  L.~Silvestrini, \textit{Nucl. Phys.} B\textbf{477}, 321 (1996).

  \item \phantomsection\label{ch28-ref-13}
  M.~K.~Gaillard and B.~W.~Lee, \textit{Phys. Rev.} D\textbf{10}, 897
  (1974).

  \item \phantomsection\label{ch28-ref-14}
  J.~Ellis and D.~V.~Nanopoulos, Ref.~\hyperref[ch28-ref-12]{12}.
  Detailed results are given by F.~Gabbiani and A.~Masiero,
  \textit{Nucl. Phys.} B\textbf{322}, 235 (1989); J.~S.~Hagelin,
  S.~Kelley, and T.~Tanaka, \textit{Nucl. Phys.} B\textbf{415}, 293
  (1994). The most complete treatment is by D.~Sutter, Stanford
  University Ph.~D. thesis (unpublished) and S.~Dimopoulos and
  D.~Sutter, Ref.~\hyperref[ch28-ref-11]{11}.

  \item[14a.] \phantomsection\label{ch28-ref-14a}
  M.~Dine, R.~Leigh, and A.~Kagan, \textit{Phys. Rev.} D\textbf{48},
  4269 (1993).

  \item \phantomsection\label{ch28-ref-15}
  For more recent reviews, see Y.~Grossman, Y.~Nir, and R.~Rattazzi,
  in \textit{Heavy Flavours II}, A.~J.~Buras and M.~Lindner, eds.
  (World Scientific, Singapore, 1998); A.~Masiero and L.~Silvestrini,
  in \textit{Perspectives on Supersymmetry},
  Ref.~\hyperref[ch28-ref-10]{10}.

  \item \phantomsection\label{ch28-ref-16}
  J.~Ellis and M.~K.~Gaillard, \textit{Nucl. Phys.} B\textbf{150}, 141
  (1979); D.~V.~Nanopoulos, A.~Yildiz, and P.~H.~Cox,
  \textit{Ann. Phys. (N.Y.)} \textbf{127}, 126 (1980); M.~B.~Gavela,
  A.~Le Yaouanc, L.~Oliver, O.~P\`ene, J.-C.~Raynal, and T.~N.~Pham,
  \textit{Phys. Lett.} \textbf{109B}, 215 (1982); B.~H.~J.~McKellar,
  S.~R.~Choudhury, X-G.~He, and S.~Pakvasa, \textit{Phys. Lett.}
  B\textbf{197}, 556 (1987).

  \item[16a.] \phantomsection\label{ch28-ref-16a}
  P.~G.~Harris \textit{et al.}, \textit{Phys. Rev. Lett.} \textbf{82},
  904 (1999).

  \item \phantomsection\label{ch28-ref-17}
  J.~Ellis, S.~Ferrara, and D.~V.~Nanopoulos, \textit{Phys. Lett.}
  \textbf{114B}, 231 (1982); J.~Polchinski and M.~B.~Wise,
  \textit{Phys. Lett.} \textbf{125B}, 393 (1983); M.~Dugan,
  B.~Grinstein, and L.~Hall, \textit{Nucl. Phys.} B\textbf{255}, 413
  (1985).

  \item \phantomsection\label{ch28-ref-18}
  R.~Arnowitt, J.~Lopez, and D.~Nanopoulos, \textit{Phys. Rev.}
  D\textbf{42}, 2423 (1990); R.~Arnowitt, M.~Duff, and K.~Stelle,
  \textit{Phys. Rev.} D\textbf{43}, 3085 (1991); Y.~Kizuri and
  N.~Oshimo, \textit{Phys. Rev.} D\textbf{45}, 1806 (1992).

  \item \phantomsection\label{ch28-ref-19}
  S.~Weinberg, \textit{Phys. Rev. Lett.} \textbf{63}, 2333 (1989);
  D.~Dicus, \textit{Phys. Rev.} D\textbf{41}, 999 (1990); J.~Dai,
  H.~Dykstra, R.~G.~Leigh, S.~Paban, and D.~A.~Dicus,
  \textit{Phys. Lett.} B\textbf{237}, 216 (1990); E.~Braaten,
  C.~S.~Li, and T.~C.~Yuan, \textit{Phys. Rev. Lett.} \textbf{64},
  1709 (1990); A.~De R\'ujula, M.~B.~Gavela, O.~Pene, and F.~J.~Vegas,
  \textit{Phys. Lett.} B\textbf{245}, 640 (1990); R.~Arnowitt,
  M.~J.~Duff, and K.~S.~Stelle, Ref.~\hyperref[ch28-ref-18]{18};
  T.~Ibrahim and P.~Nath, \textit{Phys. Lett.} \textbf{148B}, 98
  (1998).

  \item \phantomsection\label{ch28-ref-20}
  K.~S.~Babu, C.~Kolda, J.~March-Russell, and F.~Wilczek,
  \textit{Phys. Rev.} D\textbf{59}, 016004 (1999).

  \item \phantomsection\label{ch28-ref-21}
  R.~Arnowitt, M.~J.~Duff, and K.~S.~Stelle,
  Ref.~\hyperref[ch28-ref-18]{18}. In terms of the functions $J_1$ and
  $J_2$ in this reference, the function $J$ here is taken as
  $2J_1+\frac{2}{3}J_2$, under the assumption that the graphs with the
  gluon attached to the squark or gluino lines add constructively.

  \item \phantomsection\label{ch28-ref-22}
  H.~Georgi and L.~Randall, \textit{Nucl. Phys.} B\textbf{276}, 241
  (1980); A.~Manohar and H.~Georgi, \textit{Nucl. Phys.} B\textbf{238},
  189 (1984); S.~Weinberg, Ref.~\hyperref[ch28-ref-19]{19}.

  \item \phantomsection\label{ch28-ref-23}
  W.~Fischler, S.~Paban, and S.~Thomas, \textit{Phys. Lett.}
  B\textbf{289}, 373 (1992).

  \item \phantomsection\label{ch28-ref-24}
  J.~Ellis and D.~V.~Nanopoulos, Ref.~\hyperref[ch28-ref-12]{12};
  F.~Gabbiani and A.~Masiero, Ref.~\hyperref[ch28-ref-14]{14};
  F.~Dine, A.~Kagan, and S.~Samuel, \textit{Phys. Lett.}
  B\textbf{243} 250 (1990); F.~Gabbiani, E.~Gabrielli, A.~Masiero,
  and L.~Silvestrini, Ref.~\hyperref[ch28-ref-12]{12}.

  \item \phantomsection\label{ch28-ref-25}
  S.~Weinberg, \textit{Phys. Rev. Lett.} \textbf{40}, 223 (1978);
  F.~Wilczek, \textit{Phys. Rev. Lett.} \textbf{40}, 279 (1978).

  \item \phantomsection\label{ch28-ref-26}
  A.~Brignole, J.~Ellis, G.~Ridolfi, and F.~Zwirner,
  \textit{Phys. Lett.} B\textbf{271}, 123 (1991); M.~Carena,
  M.~Quiros, and C.~E.~M.~Wagner, \textit{Nucl. Phys.} B\textbf{461},
  407 (1996); S.~Heinemayer, W.~Hollik, and G.~Weiglein,
  hep-ph/9812472, hep-ph/9903404, hep-ph/9903504, to be published. The
  numerical results for $m_h$ cited here are taken from calculations
  quoted by S.~Dawson, Ref.~\hyperref[ch28-ref-10]{10}.

  \item \phantomsection\label{ch28-ref-27}
  R.~Barate \textit{et al.} (ALEPH collaboration),
  \textit{Phys. Lett.} B\textbf{412}, 173 (1997).

  \item[27a.] \phantomsection\label{ch28-ref-27a}
  M.~Gr\"unewald and D.~Karlen, in \textit{Proceedings of the XXIX
  International Conference on High Energy Nuclear Physics},
  A.~Astbury, D.~Axen, and J.~Robinson, eds. (TRIUMF, Vancouver, 1999).

  \item \phantomsection\label{ch28-ref-28}
  R.~Barate \textit{et al.} (ALEPH collaboration), 1999 CERN preprint
  EP-99-011, to be published in \textit{Phys. Lett.} A lower bound of
  54.5 GeV had been obtained earlier from experiments on
  $e^+e^-$ annihilation at 130 to 172 GeV by P.~Abreu
  \textit{et al.} (DELPHI collaboration), \textit{Phys. Lett.}
  B\textbf{420}, 140 (1998).

  \item \phantomsection\label{ch28-ref-29}
  A.~J.~Buras, M.~Misiak, M.~M\"unz, and S.~Pokorski,
  \textit{Nucl. Phys.} B\textbf{424}, 374 (1994).

  \item[29a.] \phantomsection\label{ch28-ref-29a}
  R.~Barate \textit{et al.} (ALEPH collaboration), 1999 CERN preprint
  EP-99-014, to be published in \textit{Eur. Phys. J.}

  \item \phantomsection\label{ch28-ref-30}
  M.~Dine, W.~Fischler, and M.~Srednicki, \textit{Nucl. Phys.}
  B\textbf{189}, 575 (1981); S.~Dimopoulos and S.~Raby,
  \textit{Nucl. Phys.} B\textbf{192}, 353 (1982); M.~Dine and
  W.~Fischler, \textit{Phys. Lett.} \textbf{110B}, 227 (1982);
  \textit{Nucl. Phys.} B\textbf{204}, 346 (1982); C.~Nappi and
  B.~Ovrut, \textit{Phys. Lett.} \textbf{113B}, 175 (1982);
  L.~Alvarez-Gaum\'e, M.~Claudson, and M.~Wise, \textit{Nucl. Phys.}
  B\textbf{207}, 96 (1982); S.~Dimopoulos and S.~Raby,
  \textit{Nucl. Phys.} B\textbf{219}, 479 (1983). This class of models
  was revived by M.~Dine and A.~E.~Nelson, \textit{Phys. Rev.}
  D\textbf{48}, 1277 (1993); D\textbf{51}, 1362 (1995); J.~Bagger,
  Ref.~\hyperref[ch28-ref-10]{10}; M.~Dine, A.~Nelson, and Y.~Shirman,
  \textit{Phys. Rev.} D\textbf{51}, 1362 (1995); M.~Dine, A.~Nelson,
  Y.~Nir, and Y.~Shirman, \textit{Phys. Rev.} D\textbf{53}, 2658
  (1996). For reviews, see C.~Kolda, \textit{Nucl. Phys. Proc. Suppl.}
  \textbf{62}, 266 (1998); G.~F.~Giudice and R.~Rattazzi,
  hep-ph/9801271, to be published in \textit{Phys. Rep}; S.~L.~Dubovsky,
  D.~S.~Gorbunov, and S.~V.~Troitsky, hep-ph/9905466, to be published.
  The phenomenology of these models is described by S.~Dimopoulos,
  S.~Thomas, and J.~D.~Wells, \textit{Nucl. Phys.} B\textbf{488}, 39
  (1997).

  \item \phantomsection\label{ch28-ref-31}
  S.~Dimopoulos, G.~F.~Giudice, and A.~Pomerol, \textit{Phys. Lett.}
  \textbf{389B}, 37 (1997); S.~P.~Martin, \textit{Phys. Rev.}
  D\textbf{55}, 3177 (1997).

  \item \phantomsection\label{ch28-ref-32}
  G.~F.~Giudice and R.~Rattazi, \textit{Nucl. Phys.} B\textbf{511}, 25
  (1998). This work has been extended by N.~Arkani-Hamed,
  G.~F.~Giudice, M.~A.~Luty, and R.~Rattazzi, \textit{Phys. Rev.}
  D\textbf{58}, 115005 (1998).

  \item \phantomsection\label{ch28-ref-33}
  N.~Seiberg, \textit{Phys. Lett.} B\textbf{318}, 469 (1993).

  \item \phantomsection\label{ch28-ref-34}
  K.~S.~Babu, C.~Kolda, and F.~Wilczek, \textit{Phys. Rev. Lett.}
  \textbf{77}, 3070 (1996).

  \item \phantomsection\label{ch28-ref-35}
  J.~E.~Kim and H-P.~Nilles, \textit{Phys. Lett.} \textbf{138B}, 150
  (1984); J.~Ellis, J.~F.~Gunion, H.~E.~Haber, L.~Roszkowski, and
  F.~Zwirner, \textit{Phys. Rev.} D\textbf{39}, 844 (1989); E.~J.~Chun,
  J.~E.~Kim, and H-P.~Nilles, \textit{Nucl. Phys.} B\textbf{370}, 105
  (1992); M.~Dine and A.~E.~Nelson,
  Ref.~\hyperref[ch28-ref-30]{30}; M.~Dine, A.~E.~Nelson, Y.~Nir, and
  Y.~Shirman, Ref.~\hyperref[ch28-ref-30]{30}; G.~Dvali, G.~F.~Giudice,
  and A.~Pomerol, \textit{Nucl. Phys.} B\textbf{478}, 31 (1996);
  S.~Dimopoulos, G.~Dvali, and R.~Rattazzi, \textit{Phys. Lett.}
  \textbf{413B}, 336 (1997); H-P.~Nilles and N.~Polonsky,
  \textit{Nucl. Phys.} B\textbf{484}, 33 (1997); G.~Cleaver,
  M.~Cveti\v{c}, J.~R.~Espinosa, L.~Everett, and P.~Langacker,
  \textit{Phys. Rev.} D\textbf{57}, 2701 (1998); P.~Langacker,
  N.~Polonsky, and J.~Wang, hep-ph/9905252, to be published; J.~E.~Kim,
  hep-ph/9901204, to be published.

  \item[35a.] \phantomsection\label{ch28-ref-35a}
  S.~Raby, \textit{Phys. Lett.} B\textbf{422}, 158 (1998).

  \item[35b.] \phantomsection\label{ch28-ref-35b}
  D.~A.~Dicus, B.~Dutta, and S.~Nandi, \textit{Phys. Rev. Lett.}
  \textbf{78}, 3055 (1997); \textit{Phys. Rev.} D\textbf{56}, 5748
  (1997).

  \item \phantomsection\label{ch28-ref-36}
  S.~Weinberg, \textit{Phys. Rev. Lett.} \textbf{43}, 1566 (1979);
  F.~Wilczek and A.~Zee, \textit{Phys. Rev. Lett.} \textbf{43}, 1571
  (1979).

  \item \phantomsection\label{ch28-ref-37}
  S.~Weinberg, \textit{Phys. Rev.} D\textbf{26}, 287 (1982); N.~Sakai
  and T.~Yanagida, \textit{Nucl. Phys.} B\textbf{197}, 533 (1982).
  These articles are reprinted in \textit{Supersymmetry},
  Ref.~\hyperref[ch28-ref-2]{2}.

  \item \phantomsection\label{ch28-ref-38}
  S.~Weinberg, Ref.~\hyperref[ch28-ref-36]{36}.

  \item \phantomsection\label{ch28-ref-39}
  J.~Ellis, J.~S.~Hagelin, D.~V.~Nanopoulos, and K.~Tamvakis,
  \textit{Phys. Lett.} \textbf{124B}, 484 (1983); V.~M.~Belyaev and
  M.~I.~Vysotsky, \textit{Phys. Lett.} \textbf{127B}, 215 (1983).

  \item \phantomsection\label{ch28-ref-40}
  V.~Lucas and S.~Raby, \textit{Phys. Rev.} D\textbf{55}, 6986 (1997).

  \item \phantomsection\label{ch28-ref-41}
  This estimate is taken from a compilation of non-supersymmetric
  calculations of the total two-body decay proton decay rate in terms
  of the superheavy gauge boson masses $M_X$ by P.~Langacker, in
  \textit{Proceedings of the 1983 Annual Meeting of the Division of
  Particles and Fields of the American Physical Society} (American
  Institute of Physics, New York, 1983): 251. To express the result in
  terms of $g_6$, I have assumed that, in effect, the coupling $g_6$
  that was used in these calculations was given by
  $g_6=g^2(M_X)/M_X^2$, where $g(M_X)$ has the value appropriate to
  non-supersymmetric theories, with
  $g^2(M_X)/4\pi\simeq 1/41$.

  \item \phantomsection\label{ch28-ref-42}
  For more detailed (mostly model-dependent) calculations, including
  renormalization group corrections to $g_5$, see S.~Dimopoulos,
  S.~Raby, and F.~Wilczek, \textit{Phys. Lett.} \textbf{112B}, 133
  (1982); J.~Ellis, D.~V.~Nanopoulos, and S.~Rudaz,
  \textit{Nucl. Phys.} B\textbf{202}, 43 (1982); W.~Lang,
  \textit{Nucl. Phys.} B\textbf{203}, 277 (1982); J.~Ellis,
  J.~S.~Hagelin, D.~V.~Nanopoulos, and K.~Tamvakis,
  Ref.~\hyperref[ch28-ref-39]{39}; V.~M.~Belyaev and M.~I.~Vysotsky,
  Ref.~\hyperref[ch28-ref-39]{39}; L.~E.~Ib\'a\~nez and C.~Mu\~noz,
  \textit{Nucl. Phys.} B\textbf{245}, 425 (1984); P.~Nath,
  A.~H.~Chamseddine, and R.~Arnowitt, \textit{Phys. Rev.} D\textbf{32},
  2385 (1985); J.~Hisano, H.~Murayama, and T.~Yanagida,
  \textit{Nucl. Phys.} B\textbf{402}, 46 (1993); V.~Lucas and S.~Raby,
  Ref.~\hyperref[ch28-ref-40]{40}. For a review, see P.~Nath and
  R.~Arnowitt, \textit{Phys. Atom. Nucl.} \textbf{61}, 975 (1997).

  \item[42a.] \phantomsection\label{ch28-ref-42a}
  M.~Takita \textit{et al.}, in \textit{Proceedings of the XXIX
  International Conference on High Energy Nuclear Physics},
  Ref.~\hyperref[ch28-ref-27a]{27a}.

  \item \phantomsection\label{ch28-ref-43}
  K.~S.~Babu, J.~C.~Pati, and F.~Wilczek, \textit{Phys. Lett.}
  \textbf{423B}, 337 (1998).

  \item \phantomsection\label{ch28-ref-44}
  V.~Lucas and S.~Raby, Ref.~\hyperref[ch28-ref-40]{40}; T.~Goto and
  T.~Nihei, \textit{Phys. Rev.} D\textbf{59}, 115009 (1999).

  \item \phantomsection\label{ch28-ref-45}
  L.~Ib\'a\~nez and G.~Ross, \textit{Nucl. Phys.} B\textbf{368}, 3
  (1991).
\end{enumerate}
